% =========================
% SkillBridge Project Report (LaTeX)
% CSI 4999: Senior Capstone Project - Winter 2026
% Group of 4
% =========================
\documentclass[11pt]{article}

\usepackage[margin=1in]{geometry}
\usepackage{graphicx}
\usepackage{array}
\usepackage{booktabs}
\usepackage{longtable}
\usepackage{hyperref}
\usepackage{enumitem}
\usepackage{xcolor}
\usepackage{float}
\usepackage{tikz}
\usepackage{caption}
\usepackage{setspace}
\usepackage{fancyhdr}

\hypersetup{
  colorlinks=true,
  linkcolor=blue,
  urlcolor=blue,
  citecolor=blue
}

\setlength{\parskip}{0.6em}
\setlength{\parindent}{0pt}

\pagestyle{fancy}
\fancyhf{}
\lhead{CSI 4999: Senior Capstone Project}
\rhead{Winter 2026}
\cfoot{\thepage}

% ---------- Convenience ----------
\newcommand{\ProjectName}{SkillBridge} % (Rubric: project name one or two words only)
\newcommand{\Course}{CSI 4999: Senior Capstone Project}
\newcommand{\Term}{Winter 2026}
\newcommand{\Dept}{Computer Science and Engineering Department}
\newcommand{\School}{School of Engineering and Computer Science}
\newcommand{\Univ}{Oakland University}
\newcommand{\Instructor}{Dr.\ Hadeel Mohammed Jawad}

% ---------- Update these ----------
\newcommand{\GroupNumber}{\#14} % <-- change to your actual group number
\newcommand{\VideoLink}{https://www.loom.com/share/249216f7cc3841a38c9b1cc6b66d704b} % <-- paste your video URL

\begin{document}

% =========================
% Cover Page  (Rubric #1)
% =========================
\begin{titlepage}
  \centering
  \vspace*{0.8cm}

  {\Large \Univ \par}
  \vspace{0.15cm}
  {\large \School \par}
  {\large \Dept \par}

  \vspace{0.9cm}

  {\Huge \bfseries \ProjectName \par}
  \vspace{0.25cm}
  {\large \itshape Portfolio Dashboard + Job Fit Analyzer \par}

  \vspace{0.75cm}

  % (Rubric #1) Logo or icon required
  \includegraphics[width=0.33\textwidth]{logo.png}

  \vfill

  {\large \Course \par}
  {\large \Term \par}
  \vspace{0.2cm}
  {\large Instructor: \Instructor \par}
  \vspace{0.2cm}
  {\large Date: \today \par}

  \vspace{0.5cm}
  {\large Group \GroupNumber \par}

  \vspace{0.3cm}
  {\large Team Members:\par}
  \vspace{0.1cm}
    Justin Elia\\
    Jennifer Gonzalez\\
    Spencer Roeren\\
    Cordell Stonecipher\\

  \vspace{0.5cm}
  % (Rubric #1) Video link required
  {\large Video Link: \href{\VideoLink}{\VideoLink}\par}

  \vspace{0.3cm}

\end{titlepage}

\tableofcontents
\newpage

% =========================
\section{Project Description} % (Rubric #2)
\ProjectName\ is a database-driven web application that helps students and early-career professionals
consolidate their skills, projects, and evidence into a single portfolio dashboard, then measure job fit
by pasting real job postings found online.

The core problem is that users often have valuable experience scattered across resumes, GitHub,
coursework, and personal notes, but no centralized system that (1) tracks skills with evidence,
(2) organizes projects and artifacts, and (3) evaluates readiness for a specific job posting.
\ProjectName\ addresses this by providing a structured portfolio dashboard and a job posting analyzer
that computes match coverage, identifies missing skills, and recommends next actions.

A key design decision is that job posting data is provided by users (pasted job descriptions).
This keeps the system simple, avoids scraping restrictions, and allows the job library to grow organically.
Over time, community-submitted job postings are aggregated into a role library with skill weights derived
from observed frequencies.

% =========================
\section{Targeted Audience} % (Rubric #3)
\begin{itemize}[leftmargin=*]
  \item \textbf{Primary:} Undergraduate students and new graduates preparing for internships, jobs, and capstone projects.
  \item \textbf{Secondary:} Career services staff and instructors who want a high-level view of skill development trends.
  \item \textbf{Tertiary:} Student organizations or bootcamps that want a lightweight skill tracking + job readiness tool.
\end{itemize}

% =========================
\section{Source of Revenue} % (Rubric #4)
\begin{itemize}[leftmargin=*]
  \item \textbf{Freemium:} Free dashboard + limited job analyses per month; paid tier for unlimited analyses and exports.
  \item \textbf{Institutional licensing:} Department/career center subscription for cohort analytics and administrative tools.
  \item \textbf{Premium exports:} Paid PDF portfolio exports and interview-prep gap reports derived from match results.
\end{itemize}

% =========================
\section{Similar Systems and Differentiation} % (Rubric #5)
\subsection{Similar Systems}
\begin{itemize}[leftmargin=*]
  \item \textbf{LinkedIn}: Stores experience/skills, but limited evidence linking and transparent job-fit scoring.
  \item \textbf{GitHub + portfolio sites}: Strong for projects, but lack structured skill tracking and job-posting comparison.
  \item \textbf{ATS resume scanners}: Compare resume to job postings, but do not maintain a long-term dashboard with evidence.
  \item \textbf{Learning platforms (badges/certificates)}: Track learning events but not a full skill-evidence graph plus job matching.
\end{itemize}

\subsection{How \ProjectName\ is Different}
\begin{itemize}[leftmargin=*]
  \item \textbf{Evidence-based skills:} Skills can be linked to projects and artifacts, reducing unverifiable claims.
  \item \textbf{User-provided postings:} No scraping required; users paste job descriptions directly for analysis and storage.
  \item \textbf{Community role library:} Role skill weights emerge from aggregated postings over time.
  \item \textbf{Granular skills:} Libraries (NumPy, Pandas, PostgreSQL) remain distinct skills rather than collapsing into parent labels.
\end{itemize}

% =========================
\section{Users (Two Types) and Feature Access} % (Rubric #9)
\textbf{User Type 1: Student (Standard User)} \\
Can manage portfolio content and run job analyses, but has no moderation or role-weight controls.

\textbf{User Type 2: Admin (Moderator/Curator)} \\
Can review community job postings, manage taxonomy/aliases, and generate role-level analytics.

\vspace{0.2cm}
\textbf{Access Difference (Example):}
\begin{itemize}[leftmargin=*]
  \item Student sees: Portfolio + Resume Import + Job Analyzer features (Subsystems 1--3).
  \item Admin sees: All student features + Admin/Library features (Subsystem 4).
\end{itemize}

% =========================
\section{Subsystems and Features (Group of 4)} % (Rubric #6, #8, #9)

\subsection{Subsystem 1: Portfolio Dashboard \& Evidence Mapping (Student)}
\begin{enumerate}[leftmargin=*]
  \item Create a project entry (add project metadata to database).
  \item Link a skill to a project (many-to-many mapping insert).
  \item Add evidence artifact (URL/file/notes) and associate with skill/project.
  \item Run query: dashboard summary (top skills, recent projects, evidence coverage).
\end{enumerate}

\subsection{Subsystem 2: Skill Management (Granular Skills) (Student)}
\begin{enumerate}[leftmargin=*]
  \item Add a skill with category/tags (insert skill).
  \item Update skill proficiency + last-used date (update skill metadata).
  \item View skill detail page with linked projects/evidence (join query).
  \item Run query: skills missing evidence or low-evidence skills (gap query).
\end{enumerate}

\subsection{Subsystem 3: Resume Import \& Confirmation (Student)}
\begin{enumerate}[leftmargin=*]
  \item Paste/upload resume text and store a resume snapshot (insert snapshot).
  \item Extract suggested skills using taxonomy + aliases (query + extraction results).
  \item Confirm/reject/edit extracted items and persist confirmed mappings (insert junction rows).
  \item Save confirmed resume-derived skills/projects into dashboard entities (insert/update records).
\end{enumerate}

\subsection{Subsystem 4: Job Library, Role Weights \& Moderation (Admin)}
\begin{enumerate}[leftmargin=*]
  \item Review and approve/reject community-submitted job postings (moderation status update).
  \item Tag a posting by role and store role metadata (insert/update role assignment).
  \item Aggregate postings by role and compute skill weights (group-by queries + derived table update).
  \item Maintain taxonomy: manage skill aliases and relationships (CRUD on alias/relation tables).
\end{enumerate}

\vspace{0.2cm}

% =========================
\section{Use Case Diagrams (Draw.io Exports)} % (Rubric #7)


\vspace{0.2cm}
\textbf{Subsystem 1 Use Case Diagram}\\
\begin{figure}[H]
  \centering
  \fbox{\includegraphics[width=0.95\linewidth]{uc_subsystem1.png}}
  \caption{Use Case Diagram: Portfolio Dashboard \& Evidence Mapping}
\end{figure}

\textbf{}Subsystem 2 Use Case Diagram}\\
\begin{figure}[H]
  \centering
  \fbox{\includegraphics[width=0.95\linewidth]{uc_subsystem2.png}}
  \caption{Use Case Diagram: Skill Management (Granular Skills)}
\end{figure}

\textbf{Subsystem 3 Use Case Diagram}\\
\begin{figure}[H]
  \centering
  \fbox{\includegraphics[width=0.95\linewidth]{uc_subsystem3.png}}
  \caption{Use Case Diagram: Resume Import \& Confirmation}
\end{figure}

\textbf{Subsystem 4 Use Case Diagram}\\
\begin{figure}[H]
  \centering
  \fbox{\includegraphics[width=0.95\linewidth]{uc_subsystem4.png}}
  \caption{Use Case Diagram: Job Library, Role Weights \& Moderation (Admin)}
\end{figure}

% =========================
\section{Database Design (Initial ERD + Tables)} % (Rubric #10)
The system uses a relational database to store users, skills, projects, evidence, resumes, job postings, and role aggregates.
More tables are included than the minimum to support the four subsystems and enable realistic queries.

\subsection{ERD (Initial Design)}
\begin{figure}[H]
  \centering
  \fbox{\includegraphics[width=0.95\linewidth]{erd.png}}
  \caption{Initial ERD (Entity Relationship Diagram)}
\end{figure}

\subsection{Core Tables (Draft)}
\begin{longtable}{p{0.26\linewidth} p{0.70\linewidth}}
\toprule
\textbf{Table} & \textbf{Purpose / Key Columns} \\
\midrule
\endhead

User & user\_id (PK), role (student/admin), created\_at \\
Skill & skill\_id (PK), name, category, description \\
SkillAlias & alias\_id (PK), skill\_id (FK), alias\_text \\
SkillRelation & parent\_skill\_id (FK), child\_skill\_id (FK), relation\_type \\
Project & project\_id (PK), user\_id (FK), title, summary, status, dates \\
ProjectSkill & project\_id (FK), skill\_id (FK) (junction) \\
Evidence & evidence\_id (PK), user\_id (FK), project\_id (FK), skill\_id (FK), type, url, notes \\
ResumeSnapshot & resume\_id (PK), user\_id (FK), raw\_text, created\_at \\
JobPosting & posting\_id (PK), submitter\_user\_id (FK), title, company, raw\_text, status, created\_at \\
PostingSkill & posting\_id (FK), skill\_id (FK), confidence, evidence\_phrase \\
Role & role\_id (PK), role\_name \\
RoleSkillWeight & role\_id (FK), skill\_id (FK), weight, last\_updated \\
JobAnalysis & analysis\_id (PK), posting\_id (FK), user\_id (FK), match\_score, created\_at \\
AnalysisResult & analysis\_id (FK), skill\_id (FK), match\_type, coverage \\
\bottomrule
\end{longtable}

\subsection{Data Volume Plan (Rubric Note)}
\begin{itemize}[leftmargin=*]
  \item Main table (\texttt{JobPosting} or \texttt{JobAnalysis}): plan for \textbf{30+ rows} using pasted postings from classmates and curated examples.
  \item Supporting tables (\texttt{Skill}, \texttt{Project}, \texttt{Evidence}, etc.): plan for \textbf{10+ rows} each per user cohort.
\end{itemize}

% =========================
\section{User Interface Wireframes / Storyboard} % (Rubric #11)

\subsection{Key Screens}
\begin{enumerate}[leftmargin=*]
  \item \textbf{Dashboard:} portfolio summary, recent projects, evidence coverage.
  \item \textbf{Skills:} skill list, details, evidence gaps.
  \item \textbf{Resume Import \& Confirm:} paste/upload resume, extracted items, approve edits.
  \item \textbf{Paste Job \& Analyze Match:} paste job posting, match breakdown, score.
  \item \textbf{Admin Job Library \& Role Weights:} role table, weights, moderation controls.
\end{enumerate}

\subsection{Wireframe Images}
\begin{figure}[H]
  \centering
  \fbox{\includegraphics[width=0.95\linewidth]{ui_dashboard.png}}
  \caption{Dashboard (Wireframe)}
\end{figure}

\begin{figure}[H]
  \centering
  \fbox{\includegraphics[width=0.95\linewidth]{ui_skills.png}}
  \caption{Skills (Wireframe)}
\end{figure}

\begin{figure}[H]
  \centering
  \fbox{\includegraphics[width=0.95\linewidth]{ui_resume_import.png}}
  \caption{Resume Import \& Confirm (Wireframe)}
\end{figure}

\begin{figure}[H]
  \centering
  \fbox{\includegraphics[width=0.95\linewidth]{ui_job_analyzer.png}}
  \caption{Paste Job \& Analyze Match (Wireframe)}
\end{figure}

% =========================
\section{Responsive Design Plan} % (Rubric #12)
\begin{figure}[H]
  \centering
  \fbox{\includegraphics[width=0.95\linewidth]{mobile-web_view.png}}
  \caption{Mobile/Web Views}
\end{figure}

\subsection{Desktop View}
\begin{itemize}[leftmargin=*]
  \item Multi-column dashboard layout (skills summary, projects, recommendations).
  \item Side-by-side job posting and match breakdown panels.
  \item Admin tables with sortable columns for roles and weights.
\end{itemize}

\subsection{Mobile View}
\begin{itemize}[leftmargin=*]
  \item Single-column stacked cards with collapsible sections.
  \item Match results displayed as expandable lists (Exact/Related/Missing).
  \item Sticky action buttons for core actions (Analyze, Save, Add Skill).
\end{itemize}

% =========================
\section{System Vision Document (SVD)} % (Rubric #13)
\begin{table}[H]
\centering
\renewcommand{\arraystretch}{1.2}
\begin{tabular}{p{0.24\linewidth} p{0.68\linewidth}}
\toprule
\textbf{Element} & \textbf{Description} \\
\midrule
Vision & Centralize skills/projects into a dashboard and evaluate job fit by analyzing user-pasted job postings. \\
Target Users & Students (standard users) and Admin moderators for taxonomy, moderation, and role analytics. \\
Business Need & Reduce friction in job preparation by showing readiness, evidence-backed skills, and actionable gaps. \\
Capabilities & Portfolio dashboard, skill-evidence mapping, resume import/confirm, job analyzer, admin job library and role weights. \\
Benefits & Clear readiness metrics, better portfolio organization, more targeted skill building, evolving role expectations. \\
\bottomrule
\end{tabular}
\caption{System Vision Document (SVD) Summary}
\end{table}

% =========================
\section{Work Breakdown Structure (WBS) and Iteration Schedule} % (Rubric #14)
\subsection{WBS (High-Level)}
\begin{enumerate}[leftmargin=*]
  \item Requirements \& Planning
  \item Database \& ERD Finalization
  \item Backend APIs (portfolio + job posting ingestion)
  \item Frontend Wireframes to Working UI
  \item Resume Import + Confirmation Flow
  \item Job Posting Analyzer + Match Engine
  \item Admin Library + Role Weight Aggregation
  \item Testing + Demo Data
  \item Deployment + Video + Presentation Assets
\end{enumerate}

\subsection{Iteration Schedule (Work Sequence)}
\begin{table}[H]
\centering
\renewcommand{\arraystretch}{1.2}
\begin{tabular}{p{0.16\linewidth} p{0.76\linewidth}}
\toprule
\textbf{Iteration} & \textbf{Deliverables} \\
\midrule
Iter 1 & Finalize scope, ERD draft, database schema, initial CRUD for skills/projects/evidence. \\
Iter 2 & Implement resume import + confirmation; populate demo skills/projects/evidence. \\
Iter 3 & Implement paste job analyzer + matching; save analysis results; basic UI views. \\
Iter 4 & Admin moderation + role library aggregation + role-weight views; refine queries. \\
Iter 5 & UI polish + responsive layout, testing, deployment, finalize documentation + video. \\
\bottomrule
\end{tabular}
\caption{Iteration Schedule}
\end{table}

% =========================
\section{ADDIE Model Integration} % (Rubric #15)
\begin{itemize}[leftmargin=*]
  \item \textbf{Analyze:} Identify user pain points (scattered experience, unclear job readiness, lack of evidence mapping).
  \item \textbf{Design:} Define subsystems, user roles, database schema, and wireframes for key workflows.
  \item \textbf{Develop:} Implement database tables, APIs, matching engine, admin tools, and responsive UI components.
  \item \textbf{Implement:} Deploy application, load demo postings and portfolio artifacts, run pilot usage with classmates.
  \item \textbf{Evaluate:} Collect feedback on usability and match relevance; refine taxonomy/aliases and improve UI clarity.
\end{itemize}

% =========================
\section{Conclusion}
\ProjectName\ is a practical, database-driven web application that centralizes evidence-backed skills and projects
into a portfolio dashboard and evaluates job fit using user-pasted job postings. The four-subsystem architecture
supports two distinct user types (Student and Admin), includes 16 database-backed features excluding account actions,
and provides clear artifacts required for capstone planning: use case diagrams, ERD, wireframes, responsive design plan,
SVD, WBS, and iteration schedule.

\end{document}
